\subsection{残る四型と再開する具体的な入口}
\label{ns22:support-agenda}\label{ns20:support-agenda}

再開時は最初に \nolinkurl{00_CONTINUATION_POLICY_ja.txt} を読む．
現在の高分母支持の必要候補は\emph{四型・五十六標識}である．
全候補は五頂点を使い，二点連結，最小次数二以上，八高辺以上，$b_1\ge4$ となる．
二価頂点は高々一つであり，ある場合は $G_{5,15}$，その高近傍はともに四価である．
十本の異なる branches のうち三本以上が分母一なら，全分母一となる．
補 graph の孤立頂点を省略すると，候補表は次の通りである．
\begin{center}\small
\begin{tabular}{c|c|c|c|r}
型&高辺数&次数列&低辺補 graph&標識数\\\hline
$G_{5,14}$&8&$(43333)$&$2K_2$&15\\
$G_{5,15}$&8&$(44332)$&$P_3$&30\\
$G_{5,16}$&9&$(44433)$&$K_2$&10\\
$G_{5,17}$&10&$(44444)$&$\varnothing$&1
\end{tabular}
\end{center}
識別子を詰め直さない．全 $1024$ graph の現行検算は
\nolinkurl{verify_support_agenda_v22.py} と
\nolinkurl{support_agenda_verification_v22.json} に保持する．
候補数は整数接着，atlas，幾何的非同型類の存在数ではない．

\paragraph{Book は解決済み．}
負枝 $r=2w-1,2w$ の全尺度二族と，正枝 $L=3,4$ の全尺度四族を他の支持へ使う．
Book の大きい分母や四高中心の完成は再探索しない．
証明の前段階で用いた全尺度有限証明書は再現性のため保持する．

\paragraph{優先候補：$G_{5,15}$ の非共鳴領域と負例外．}
標識を $H=K_5-\{23,24\}$，二価 core を $2$，共有高分母を $Q=q_{34}$ とする．
通常型は四中心・十四配置，負例外型は $Q=3,5$ の二 template に限られる．
通常型の共有高 $D$ 共鳴は全て排除された．
正 reverse 側の五つの固定低データ問題
\[
((-1,-4,4),5),\quad((-1,-1,1),-9),\quad
((-1,-2,2),2),\quad((-1,-2,2),-7),\quad((-1,-2,2),9)
\]
では $T\ge2$ に限ってよい．非共鳴一次式の解を反対側の pair equation へ代入し，
正側の二次復元と合わせて，無上界の $Q$ に対する障害または構成族を求める．
負例外では負枝と正 $L=3$ の全尺度族を二高中心へ代入し，七・十三の共鳴も保持する．
共有 $Q$ が小さいことだけから中心の有限性を結論しない．

\paragraph{並行候補：三高 $D$ の二変数問題．}
固定 $(s,t)$ に対する有限復元を用いる．
$F\le st+(2(s+t-1)+1)(s-x)-2$ により $x,b,z$ が有限となり，
$y$ は元の残余式の非負整数根として全て求まる．
対称 sector $s=t,x=y$ は無上界で排除済みである．
$s,t\le12$ の全 $144$ 低行では他の全係数・全分母を尽くして模型零だったが，
$(s,t)$ 自体が全域を動く非対称の場合は未解決である．
次は判別式の平方性，単一合同式，元の pair equations の関係を調べる．
反例が見つかれば，全原方程式と正の weights を復元してその構造へ転換する．

\paragraph{他の支持と幾何への転換．}
$G_{5,14},G_{5,16},G_{5,17}$ は引き続き必要候補である．
四高側へ三高 matching 定理をそのまま適用しない．
$K_4$ を部分 graph として含むことだけでは支持全体が $K_4$ の障害を移せない．
素数支持の合同条件，特に dense mod 五の二模型の少なくとも二本の付値一を保持する．
整数候補が残れば全 core の weights，basin areas，integer framings，作用切片へ進む．
幾何的同値，不変量，別の実際の退化，open-corridor surgery への転換も，
研究全体の面白さとインパクトに応じて判断する．
